% Source: https://artofproblemsolving.com/community/c3h2703888_hello_a_problem
% Stable identity: aops-topic-2703888-q01
Let $\alpha=\dfrac{180^\circ}{7}$. The vertices lie on a circle and divide it into seven equal arcs of $2\alpha$ each. An angle at a vertex between two chords is an inscribed angle, so it is a multiple of $\alpha$. An angle at a crossing of two diagonals is half the sum of two intercepted arcs, so it is also a multiple of $\alpha$. Hence every region has all its angles in multiples of $\alpha$.

\begin{asy}
usepackage("amsmath");
size(4cm);
defaultpen(fontsize(10pt));
pair[] V;for(int k=0;k<7;++k)V[k]=dir(90+360*k/7);
for(int i=0;i<7;++i)for(int j=i+1;j<7;++j)draw(V[i]--V[j],linewidth(0.5));
for(int k=0;k<7;++k)draw(V[k]--V[(k+1)%7],linewidth(1));
\end{asy}

No three diagonals of a regular heptagon are concurrent, so there are $\dbinom74+\dbinom62=50$ regions. Of these, $35$ are triangles, forming five rotation classes of $7$. Their angle types, in units of $\alpha$, are $(1,1,5)$, $(1,2,4)$ (two classes), $(1,3,3)$ and $(2,2,3)$. All classes except $(2,2,3)$ contain an angle $\alpha$, so the number of tringles is $4\cdot7=\boxed{28}$.
